\section*{Bab \ref{ch:g8:equations} --- Perhitungan Aljabar dan Persamaan}

\begin{solution}{exo:g8:equations:1}
$5(2x + 3) = 10x + 15$.

$-3(4x - 1) = -12x + 3$.

$2(3x - 2) + 4(x + 5) = 6x - 4 + 4x + 20 = 10x + 16$.
\end{solution}

\begin{solution}{exo:g8:equations:2}
$(x+2)(x+7) = x^2 + 7x + 2x + 14 = x^2 + 9x + 14$.

$(3x+1)(2x-5) = 6x^2 - 15x + 2x - 5 = 6x^2 - 13x - 5$.

$(x-4)(x-6) = x^2 - 6x - 4x + 24 = x^2 - 10x + 24$.
\end{solution}

\begin{solution}{exo:g8:equations:3}
$4 \times 3 - 5 = 7$: ya. \quad
$2 \times 3 + 1 = 7$ dan $3 + 4 = 7$: ya. \quad
$3^2 = 9$ dan $6 \times 3 - 9 = 9$: ya --- $x = 3$ menyelesaikan
ketiganya.
\end{solution}

\begin{solution}{exo:g8:equations:4}
$x + 9 = 4$: $x = -5$ (periksa: $-5 + 9 = 4$).

$5x = -35$: $x = -7$.

$\frac x3 = 8$: $x = 24$.

$2x - 7 = 13$: $2x = 20$, $x = 10$ (periksa: $20 - 7 = 13$).
\end{solution}

\begin{solution}{exo:g8:equations:5}
$6x + 5 = 2x + 25$: $4x = 20$, jadi $x = 5$.

$9 - 4x = x - 11$: $20 = 5x$, jadi $x = 4$.

$3(x + 2) = 2x + 11$: $3x + 6 = 2x + 11$, jadi $x = 5$.
\end{solution}

\begin{solution}{exo:g8:equations:6}
$8x - 12 = 5x + 30$, jadi $3x = 42$ dan $x = 14$. Periksa:
$4(28 - 3) = 100$ dan $5(14 + 6) = 100$.
\end{solution}

\begin{solution}{exo:g8:equations:7}
``$3x - 8 = 19$'': $3x = 27$, $x = 9$.

``$2x + 5 = x + 12$'': $x = 7$.
\end{solution}

\begin{solution}{exo:g8:equations:8}
Lebarnya $x$, panjangnya $x + 5$; \omterm{def:g6:measure:perimeter}{kelilingnya}
$2\bigl(x + (x + 5)\bigr) = 4x + 10 = 58$, jadi $4x = 48$ dan
$x = 12$: lebarnya $12$ cm, panjangnya $17$ cm (periksa:
$2 \times (12 + 17) = 58$).
\end{solution}

\begin{solution}{exo:g8:equations:9}
Bilangannya adalah $n - 1$, $n$, $n + 1$ dengan \omterm{def:g1:addition:def}{jumlah} $3n = 141$,
jadi $n = 47$: bilangannya adalah $46$, $47$, $48$.
\end{solution}

\begin{solution}{exo:g8:equations:10}
$3x + 4 > 19$: $3x > 15$, jadi $x > 5$.

$7 - 2x \geq 1$: $-2x \geq -6$, membaginya dengan $-2$ membalikkannya:
$x \leq 3$.

$5(x-1) < 2x + 7$: $5x - 5 < 2x + 7$, jadi $3x < 12$ dan $x < 4$.
\end{solution}

\begin{solution}{exo:g8:equations:11}
Sesudah $n$ bulan, Mia punya $140 + 15n$ dan Tono punya
$200 + 10n$. Mia unggul ketika
\[
140 + 15n > 200 + 10n
\ \Longleftrightarrow\
5n > 60
\ \Longleftrightarrow\
n > 12 :
\]
mulai dari bulan yang ke-$13$.
\end{solution}

\begin{solution}{exo:g8:equations:12}
Kalikan kedua ruasnya dengan $20$:
\[
5(x + 3) = 4(2x - 1)
\ \Longrightarrow\
5x + 15 = 8x - 4
\ \Longrightarrow\
19 = 3x
\ \Longrightarrow\
x = \frac{19}{3}.
\]
Periksa: $\frac{19/3 + 3}{4} = \frac{28/3}{4} = \frac{28}{12} =
\frac73$, dan $\frac{2 \times 19/3 - 1}{5} = \frac{35/3}{5} =
\frac{35}{15} = \frac73$. Sama: penyelesaiannya $\frac{19}{3}$.
\end{solution}

\begin{solution}{pb:g8:equations:1}
\textbf{1.} Sifat distributif ganda
(\cref{thm:g8:equations:expand}) dengan $c = a$, $d = b$:
\[
(a + b)(a + b) = a^2 + ab + ba + b^2 = a^2 + 2ab + b^2 ,
\]
karena $ab$ dan $ba$ adalah \omterm{def:g2:mult:def}{hasil kali} yang sama.

\textbf{2.} Penjabaran yang sama, dengan memperhatikan tandanya:
\[
(a - b)(a - b) = a^2 - ab - ba + b^2 = a^2 - 2ab + b^2 ,
\]
\[
(a + b)(a - b) = a^2 - ab + ba - b^2 = a^2 - b^2 :
\]
kali ini suku silangnya $-ab$ dan $+ba$ saling meniadakan, bukan
berlipat dua.

\textbf{3.} Persegi bersisi $a + b$ terpotong menjadi empat keping:
persegi $a \times a$ (yaitu $a^2$), persegi $b \times b$ (yaitu
$b^2$), dan \emph{dua} persegi panjang $a \times b$ (bersama-sama
menjadi $2ab$). Kesamaan $(a+b)^2 = a^2 + 2ab + b^2$ menyatakan bahwa
keempat luasnya memenuhi persegi besarnya --- itulah gambar sifat
distributif ganda dengan kedua sisinya dibelah dengan cara yang sama.

\textbf{4.} Dengan kesamaannya:
\begin{align*}
101^2 &= (100 + 1)^2 = 10\,000 + 2 \times 100 + 1 = 10\,201, \\
99^2 &= (100 - 1)^2 = 10\,000 - 200 + 1 = 9\,801, \\
49 \times 51 &= (50 - 1)(50 + 1) = 50^2 - 1^2 = 2\,499 .
\end{align*}

\textbf{5.} Untuk $a = 3$, $b = 4$: $(3 + 4)^2 = 49$ tetapi
$3^2 + 4^2 = 9 + 16 = 25$. Rumus Zahra melupakan kedua persegi
panjang $a \times b$ pada gambarnya --- yaitu $2ab = 24$ yang hilang,
dan memang $25 + 24 = 49$.

\textbf{6.} $1$; $1 + 3 = 4$; $1 + 3 + 5 = 9$;
$1 + 3 + 5 + 7 = 16$; $1 + 3 + 5 + 7 + 9 = 25$. Hasilnya adalah
kuadrat sempurna $1^2, 2^2, 3^2, 4^2, 5^2$.

\textbf{7.} \omterm{def:g3:numbers:evenodd}{Bilangan ganjilnya} bermula di $1$ lalu naik dengan
langkah $2$: yang ke-$n$ tercapai sesudah $n - 1$ langkah sebesar $2$
dari $1$, jadi nilainya $1 + 2(n - 1) = 2n - 1$. Periksa:
$n = 1, 2, 3$ memberi $1, 3, 5$.

\textbf{8.} Menurut kesamaan yang pertama,
\[
(n + 1)^2 - n^2 = n^2 + 2n + 1 - n^2 = 2n + 1 :
\]
\omterm{ex:g1:subtraction:difference}{selisih} antara dua kuadrat yang berurutan persis \omterm{def:g3:numbers:evenodd}{bilangan ganjil} yang
ke-$(n+1)$ (pertanyaan 7 dengan $n + 1$ menggantikan $n$:
$2(n+1) - 1 = 2n + 1$).

\textbf{9.} Bangunlah perseginya satu demi satu. Mulailah dengan satu
sel: $1 = 1^2$. Untuk berpindah dari persegi bersisi $1$ ke persegi
bersisi $2$, tambahkan $2 \times 1 + 1 = 3$ sel (pertanyaan
8): $1 + 3 = 2^2$. Untuk berpindah ke sisi $3$, tambahkan
$2 \times 2 + 1 = 5$ sel: $1 + 3 + 5 = 3^2$. Tiap \omterm{def:g3:numbers:evenodd}{bilangan ganjil}
yang baru persis lapisan berbentuk L yang melengkapi persegi
berikutnya, jadi sesudah menambahkan $n$ \omterm{def:g3:numbers:evenodd}{bilangan ganjil} pertamanya
persegi bersisi $n$ menjadi lengkap: $1 + 3 + \dots + (2n - 1) = n^2$.

\textbf{10.} \omterm{def:g3:numbers:evenodd}{Bilangan ganjil} dari $1$ sampai $99$ adalah $50$ yang
pertama ($2n - 1 = 99$ memberi $n = 50$), jadi
\[
1 + 3 + 5 + \dots + 99 = 50^2 = 2\,500 .
\]
Dan $1 + 3 + \dots + (2n - 1) = 400 = 20^2$ menuntut $n = 20$:
\emph{dua puluh} \omterm{def:g3:numbers:evenodd}{bilangan ganjil} pertamanya ($1$ sampai $39$)
berjumlah $400$.

\textbf{11.} $x^2 + 6x + 9 = x^2 + 2 \times 3 \times x + 3^2 =
(x + 3)^2$: kesamaan pertama dengan $a = x$, $b = 3$. Jadi
\omterm{def:g8:equations:def}{persamaannya} berbunyi $(x + 3)^2 = 0$. Kuadrat bernilai \omterm{def:g1:counting:zero}{nol} hanya
ketika bilangannya sendiri \omterm{def:g1:counting:zero}{nol} (\cref{thm:g8:negprod:rules}: bilangan
yang bukan \omterm{def:g1:counting:zero}{nol} punya jarak yang bukan \omterm{def:g1:counting:zero}{nol}, jadi kuadratnya pun punya
jarak yang bukan \omterm{def:g1:counting:zero}{nol}). Karena itu $x + 3 = 0$: satu-satunya
penyelesaiannya adalah $x = -3$.

\textbf{12.} Menurut \cref{thm:g8:negprod:rules}, jarak sebuah
\omterm{def:g2:mult:def}{hasil kali} adalah \omterm{def:g2:mult:def}{hasil kali} kedua jaraknya. Kalau kedua faktornya
bukan \omterm{def:g1:counting:zero}{nol}, kedua jaraknya bukan \omterm{def:g1:counting:zero}{nol}, jadi jarak \omterm{def:g2:mult:def}{hasil kalinya} pun
bukan \omterm{def:g1:counting:zero}{nol}: \omterm{def:g2:mult:def}{hasil kalinya} tidak dapat bernilai $0$. Karena itu \omterm{def:g2:mult:def}{hasil
kali} yang \omterm{def:g1:counting:zero}{nol} memaksa sedikitnya satu faktornya \omterm{def:g1:counting:zero}{nol}. Untuk
$(x - 2)(x + 5) = 0$: entah $x - 2 = 0$ entah $x + 5 = 0$, yang
memberi dua penyelesaian $x = 2$ dan $x = -5$.

\textbf{13.} $x^2 = 49$ menjadi $x^2 - 49 = 0$, dan kesamaan yang
ketiga memfaktorkannya:
\[
x^2 - 49 = x^2 - 7^2 = (x - 7)(x + 7) = 0 .
\]
Menurut pertanyaan 12, $x = 7$ atau $x = -7$: \emph{dua}
penyelesaian. Setiap \omterm{def:g8:equations:def}{persamaan} yang diselesaikan sejauh ini pada bab
ini berderajat satu dan punya tepat satu penyelesaian; kuadrat di
dalam \omterm{def:g8:equations:def}{persamaan} dapat menghasilkan dua.

\textbf{14.} Bawalah semuanya ke ruas kiri:
$x^2 - 6x + 9 = 0$, lalu kenalilah kesamaan yang kedua:
$x^2 - 2 \times 3 \times x + 3^2 = (x - 3)^2$. Jadi \omterm{def:g8:equations:def}{persamaannya}
adalah $(x - 3)^2 = 0$, yang memaksa $x - 3 = 0$ (seperti pada
pertanyaan 11): $x = 3$ adalah satu-satunya penyelesaiannya ---
pemeriksaan \cref{exo:g8:equations:3} menemukannya, dan
pemfaktorannya membuktikan bahwa tidak ada yang lain.

\textbf{15.} Menurut kesamaan yang ketiga,
\[
2026^2 - 2025^2 = (2026 - 2025)(2026 + 2025)
= 1 \times 4051 = 4\,051 .
\]
Secara umum $(n+1)^2 - n^2 = (n + 1 - n)(n + 1 + n) = 2n + 1$:
\omterm{ex:g1:subtraction:difference}{selisih} dua kuadrat yang berurutan adalah \omterm{def:g1:addition:def}{jumlah} kedua bilangannya
--- yaitu $2n + 1$ yang sama seperti pada pertanyaan 8, yang kini
diperoleh dari kesamaan ketiganya, bukan dari yang pertama.

\end{solution}
